\section{Can the paired state follow the deflation?}
\label{sec:quench}

\subsection{The question, and the parameter that decides it}

Section~\ref{sec:gap} gave the gap that the paired state would have in equilibrium at each moment.
But a state can change at most at the rate of its own energy scale, which for a paired state is the
gap. By Eq.~(\ref{eq:growth}) the equilibrium gap, written $\Deq$ in this section, changes at the
rate $\gsum/\lambda$. The paired state can follow only if that rate is small compared with the gap
itself:
\begin{equation}
\Xid=\frac{\gsum}{\lambda\,\Deq}\ \ll\ 1 .
\label{eq:Xi}
\end{equation}
Early in the deflation the gap is exponentially small, and the condition fails however small the
rate is. A system that is carried through a change faster than it can respond falls out of
equilibrium, as in the formation of defects at a phase transition \cite{Kibble1976,Zurek1985}. The
question is what the paired state does then, and when it catches up.

\subsection{The model}

The question is answered with the standard reduced model of a superfluid, whose dynamics without
collisions is known in detail \cite{BarankovLevitovSpivak2004}. Pairs of states of opposite momenta,
with the single-particle energy $\xi_k$ measured from the Fermi energy, fill a shell
$-\Esh<\xi_k<\Esh$ with a constant density of states, and all attract one another with the
dimensionless strength $\lambda$. In the mean-field approximation each pair of states is described by
its pair correlation $\kappa$, a complex number, and its occupation $n_k$; the order parameter is
$\Delta=\lambda\int\kappa\,d\xi_k$. The equations of motion are
\begin{align}
\frac{d\kappa}{dx4}&=-i\,\big[2\,\xi_k\,\kappa-\Delta\,(1-2n_k)\big],\nonumber\\
\frac{dn_k}{dx4}&=2\,\mathrm{Im}\big(\Delta^{*}\kappa\big).
\label{eq:bloch}
\end{align}
They conserve $|2\kappa|^{2}+(1-2n_k)^{2}=1$ for every pair of states. In equilibrium
$\kappa=\Delta/2E_k$ and $n_k=\tfrac12(1-\xi_k/E_k)$ with the quasi-particle energy
$E_k=(\xi_k^{2}+\Delta^{2})^{1/2}$, and the gap equation gives
\begin{equation}
\Deq(\lambda)=\frac{\Esh}{\sinh(1/\lambda)}\ \approx\ 2\,\Esh\,e^{-1/\lambda}\quad(\lambda\ll1).
\end{equation}
The model keeps what matters for the question, a gap that depends exponentially on a coupling that
grows exponentially, and omits the dependence of the gap on the momentum; $\Esh$ stands for the
range of energies over which the attraction acts. The program starts the model in equilibrium at
$\lambda=0.15$, lets the coupling grow as $\lambda=0.15\,e^{\td}$ up to $\lambda=0.5$, and integrates
Eqs.~(\ref{eq:bloch}) for 2000 energies. Doubling the number of energies and halving the step
changes no digit of the results quoted. Besides $|\Delta|$ the program records the density $\nex$ of
\emph{broken pairs}: the number of pairs of states, per unit of $\Esh$ times the density of states,
that are not in the ground state that belongs to the order parameter of the moment. In equilibrium
it is zero.

\subsection{Results}

\begin{table*}
\caption{\label{tab:quench}The paired state while the pairing strength grows from 0.15 to 0.5, for
seven rates of deflation: the parameter $\Xid$ at the start; the smallest ratio of the order
parameter to the equilibrium gap of the moment; the clock reading at which the order parameter has
returned to within 10 per cent of the equilibrium gap, and $\Xid$ at that moment; the ratio at the
end; and the density of broken pairs at the end.}
\begin{ruledtabular}
\begin{tabular}{ccccccc}
$\gsum/\Esh$ & $\Xid$ at the start & smallest $|\Delta|/\Deq$ & $\td$ of the catch-up & $\Xid$ there & final $|\Delta|/\Deq$ & $\nex$ at the end \\
\hline
{{ROWS:tbl_quench}}
\end{tabular}
\end{ruledtabular}
\end{table*}

Table~\ref{tab:quench} and Fig.~\ref{fig:quench} show three results.

\emph{The faster the deflation, the further the order parameter falls behind.} For a rate of
$10^{-4}\,\Esh$ the order parameter drops to 0.958 of the equilibrium gap of the moment before it
recovers; for $10^{-1}\,\Esh$ it drops to 0.098.

\emph{The state catches up when $\Xid$ has fallen to a value between 0.23 and 0.52.} In six of the
seven runs the order parameter returns to within 10 per cent of the equilibrium gap at such a value.
In the fastest run it has not caught up when the coupling reaches 0.5. This confirms the estimate
(\ref{eq:Xi}) and fixes the number that ``small'' means.

\emph{The lag leaves broken pairs behind.} Their density at the end grows with the rate, from
$6.0\times10^{-6}$ to $4.1\times10^{-2}$ in the six slower runs. They are excitations of the
superfluid that the deflation itself has made; in a system without dissipation they remain.

Because $\Deq$ grows as $e^{-1/\lambda}$ while $\lambda$ grows only as $e^{\td}$, the parameter
$\Xid$ falls steeply with time. For any rate there is a moment at which $\Xid$ passes through the
window of the second result. Before it the paired state is out of equilibrium, with an order
parameter below the equilibrium value; after it the adiabatic history of Sec.~\ref{sec:gap}
applies. The slower the deflation, the earlier that moment and the smaller the gap at which it
happens: it is the moment at which the equilibrium gap has grown to between two and four times
$\gsum/\lambda$.

\begin{figure*}
\includegraphics[width=\textwidth]{fig_quench}
\caption{\label{fig:quench}The paired state while the coupling grows. (a) The order parameter
divided by the equilibrium gap of the moment, against the pairing strength, for seven rates of
deflation. (b) The smallest value of that ratio against the rate. (c) The density of broken pairs at
the end of the run against the rate.}
\end{figure*}

